% !TeX program = lualatex
% !TeX encoding = utf8
% !TeX root = ../Thermod.tex

%=========================================================
\Opensolutionfile{answer}[\currfilebase/\currfilebase-Answers]
\Writetofile{answer}{\protect\section*{\nameref*{\currfilebase}}}
\chapter{Принципи екстремумів і термодинамічної рівноваги}\label{\currfilebase}
\makeatletter
\def\input@path{{\currfilebase/}}
\makeatother
%=========================================================

% --------------------------------------------------------
\section{1}
% --------------------------------------------------------
Теорія термодинамічної рівноваги була розвинена Гіббсом на подобі теорії рівноваги механічних систем. Як відомо, в механічній системі рівноважному положенню відповідає екстремальне положення потенціальної енергії системи-мінімум, або еквівалентно-для виводу системи із положення рівноваги необхідно виконати над нею роботу. Розумно припустити, що роль потенціальної енергії в термодинаміці можуть виконувати термодинамічні потенціали $U$, $F$, $H$, $G$, що володіють, як ми встановили усіма властивостями потенціальної енергії. Щоб показати це будемо спиратись на другий закон термодинаміки, згідно з яким, кожна ізольована система еволюціонує до стану, відповідаючи максимуму ентропії. І насправді, внаслідок наявності незворотних процесів в ізольованій системі ($U = \text{соnst}$, $V = \text{соnst}$) виконується співвідношення:
\begin{equation}\label{7.1}
	dS\geq0
\end{equation}
Поки нарешті не буде досягнуто стану,при котрому подальша зміна $S$ неможлива. З цього випливає що у цьому стані $S = S_{max}$.
% --------------------------------------------------------
\section{2}
% --------------------------------------------------------
Зауважимо, що (\ref{7.1}) випливає з рівнянь термодинаміки:
\begin{equation}\label{7.2}
	dS \geq \dfrac{1}{T}(dU + PdV)
\end{equation}
Якщо припустити, що $U$ і $V$ --- соnst. Аналогічно з (\ref{7.1}) випливає, що у випадку постійних $V$ і $S$ умовою рівноваги термодинамічної системи буде мінімум внутрішньої енергії системи. $U_{min} \Rightarrow \delta U = 0,~~ (V,~S = \text{соnst})$

Насправді, з (\ref{7.1}) випливає, що при $V$ і $S$ незмінних, неврівноважені процеси протікають так, що $dU<0$, до тих пір поки не буде досягнуто мінімуму. Зауважимо, що умова максимуму ентропії і мінімуму внутрішньої енергії було виведено з умови існування незворотних процесів, для котрих справджується (\ref{7.1}). З точки зору математики ці умови слід записати таким чином:
\begin{equation}\label{7.3}
	\delta S = 0; ~~~~ \delta^2 S < 0,
\end{equation}
що гарантує опуклість функції $S=S(U,V)$ в точці $S=S_{max}$. Аналогічно для $U=U(S,V)$:
\begin{equation}\label{7.4}
	\delta U =0; ~~~~ \delta^2 U < 0,
\end{equation}
що гарантує увігнутість функції $U(S,V)$ в точці $U=U_{min}$ і тим самим стан $U_{min}$ є стійким. Зміст запису $\delta S$ і $\delta U$ маємо розуміти так, що мова йде про віртуальні зміни при збереженні умов $(dU=dV=0)$ або $(dV=dS=0)$, тобто змінах зовнішніх і внутрішніх параметрів сумісних з накладеними умовами. Наприклад, збільшення $\delta T$ в одній частині і зменшення в іншій на $\delta T$.
% --------------------------------------------------------
\section{3}
% --------------------------------------------------------
Звернемося тепер до систем що знаходяться у термостаті ($T=\text{соnst}$) при постійному об'ємі. Характеристичною функцією для такої системи буде $F=F(T,V)$. Її диференціал буде:
\begin{equation}\label{7.5}
	dF = dU -TdS
\end{equation}
Слідуючи логіці наведених міркувань неврівноважені процеси можуть лише зменшувати $P$. Дійсно, в найзагальнішому випадку:
\begin{equation}\label{7.6}
	dF\leq dU-TdS; ~~~~ dF\leq -SdT-PdV
\end{equation}
Звідси, рівноважне значення системи з $T$, $U=\text{const}$ відповідає мінімуму вільної енергії Гельмгольца:
\begin{equation}\label{7.7}
	\delta F = 0; ~~~~ \delta^2 F >0
\end{equation}

В зв'язку з тим що спад $P$ дорівнює роботі здійснюваною системою в квазістатичному процесі, то з (\ref{7.2}) випливає теорема про максимальну роботу: У процесах, температура початкового і кінцевого стану яких однакова, максимальну роботу, котру може виконати система над середовищем, можна отримати лише в оборотних процесах для яких:
\begin{equation}\label{7.8}
	A_{max} = \Delta F_{12}
\end{equation}

Ця властивість дала привід Гельмгольцу називати $P$ вільною енергією. Різниця її значень для двох станів 1 та 2 з однаковою температурою $T_0$, виражає ту частину зміни внутрішньої енергії, котра при ізотермічному оборотному процесі може піти на виконання роботи. Залишкова частина $U-F=TS$ має назву \enquote{зв'язкова енергія}. Найбільш змістовна частина теореми про максимальну роботу полягає у вимозі оборотності процесу. І насправді, розмірковуючи подібним чином отримуємо, що для адіабатно ізольованих систем $A_{max} = \Delta U$, 
що, утім, являє собою наслідок першого закону термодинаміки, але тільки для оборотних процесів. Узагальнюючи теорему, робимо висновок(укладаємо(заключаємо)), що максимальна робота може бути отримана в оборотних процесах і для розширених систем, для котрих:
\begin{equation}\label{7.9a}\tag{7.9а}
	A_{max} = \Delta G, \text{ при } T = \text{const}
\end{equation}
\begin{equation}\label{7.9b}\tag{7.9б}
	A_{max} = \Delta H, \text{ при } S = \text{const}
\end{equation}
% --------------------------------------------------------
\section{4}
% --------------------------------------------------------